%%
%% The abstract is a short summary of the work to be presented in the
%% article.
\begin{abstract}
% Agents 因其决策能力和完成复杂任务备受关注
LLM-based agents have gained considerable attention for their decision-making skills and ability to handle complex tasks.
% Motivation: 考虑到目前基于agents进行推荐的工作在多智能体合作的探索不足，且缺少开源框架，我们提出了MACRec框架。
Recognizing the current gap in leveraging agent capabilities for multi-agent collaboration in recommendation systems, we introduce \textbf{MACRec}, a novel framework designed to enhance recommendation systems through multi-agent collaboration.
% 具体介绍
Unlike existing work on using agents for user/item simulation, we aim to deploy multi-agents to tackle recommendation tasks directly.
In our framework, recommendation tasks are addressed through the collaborative efforts of various specialized agents, including \textsl{Manager}, \textsl{User/Item Analyst}, \textsl{Reflector}, \textsl{Searcher}, and \textsl{Task Interpreter}, with different working flows.
Furthermore, we provide application examples of how developers can easily use MACRec on various recommendation tasks, including rating prediction, sequential recommendation, conversational recommendation,
% interactive recommendation, click-through rates (CTR) prediction, 
and explanation generation of recommendation results.
% 开源 % The framework is available at https://github.com/wzf2000/MACRec.
% The framework and its documentation and examples are made publicly available at https://github.com/wzf2000/MACRec.
The framework and demonstration video are publicly available at \href{https://github.com/wzf2000/MACRec}{https://github.com/wzf2000/MACRec}.

% Agents 因其决策能力和完成复杂任务备受关注
% LLM-based agents have gained considerable attention for their decision-making skills and ability to handle complex tasks.
% Motivation: 考虑到目前基于agents进行推荐的工作在多智能体合作的探索不足，且缺少开源框架，我们提出了MACRec框架。
% There are various complex decision-making tasks in recommendation scenarios, so some previous studies have tried to use agents in this scenario. However, existing work focuses on using agents for user/item simulation, ignoring the multi-agent collaboration ability for recommendation. 
% In this paper, we introduce \textbf{MACRec}, a novel framework designed to enhance recommendation systems through multi-agent collaboration, to deploy multi-agents for various recommendation tasks directly.
% 具体介绍
% Unlike existing work on using agents for user/item simulation, we aim to deploy multi-agents to tackle recommendation tasks directly.
% In our framework, recommendation tasks are addressed through the collaborative efforts of various specialized agents, including \textsl{Manager}, \textsl{User/Item Analyst}, \textsl{Reflector}, \textsl{Searcher}, and \textsl{Task Interpreter}, with different working flows.
% Furthermore, we provide application examples of how developers can easily use MACRec on various recommendation tasks, including rating prediction, sequential recommendation, conversational recommendation,
% interactive recommendation, click-through rates (CTR) prediction, 
% and explanation generation of recommendation results.
% 开源 % The framework is available at https://github.com/wzf2000/MACRec.
% The framework and its documentation and examples are made publicly available at https://github.com/wzf2000/MACRec.
% The framework and demonstration video are publicly available at https://github.com/wzf2000/MACRec.

\end{abstract}


% LLM-based agents have gained considerable attention for their decision-making skills and ability to handle complex tasks. There are various complex decision-making tasks in recommendation scenarios, so some previous studies have tried to use agents in this scenario. However, existing work focuses on using agents for user/item simulation, ignoring the multi-agent collaboration ability for recommendation. In this paper, we introduce MACRec, a novel framework designed to enhance recommendation systems through multi-agent collaboration, to deploy multi-agents for various recommendation tasks directly. In our framework, recommendation tasks are addressed through the collaborative efforts of various specialized agents, including Manager, User/Item Analyst, Reflector, Searcher, and Task Interpreter, with different working flows. Furthermore, we provide application examples of how developers can easily use MACRec on various recommendation tasks, including rating prediction, sequential recommendation, conversational recommendation, and explanation generation of recommendation results. The framework and demonstration video are publicly available at https://github.com/wzf2000/MACRec.
